% !PRINTED 147
\weberpage{147}
For if $Q_0$ were intransitive in these $n-1$ figures, then by $\alpha$)

\[
\begin{aligned}
&\text{either } \Pi(n-2) \text{ divisible by } g,\\
&\text{or } \Pi(n-3)\,\Pi(2) \geqq g
\end{aligned}
\]

would have to hold; hence by (9):

\[
\begin{aligned}
&\text{either } \frac{n}{n-1} \text{ an integer, hence } n \geqq 2(n-1)\\
&\text{or } n^2 - 5n + 2 \leqq 0.
\end{aligned}
\]

The one is impossible when $n > 2$, the other when $n > 4$.

We set aside the case $n = 4$ and now consider the divisor $Q_{0,1}$ of
$Q$ which leaves the two figures 0, 1 unchanged.

Since $Q_{0,1}$, regarded as a permutation group of the $n-1$ figures, is a
transitive divisor of $Q_0$, the degree $g_1$ of $Q_{0,1}$ is, as one concludes
by applying the decomposition (8) once more, equal to $g : n-1$; hence

\weq{10}{g_1 = \frac{\Pi(n-1)}{n(n-1)} = \frac{\Pi(n-2)}{n}.}

From this we conclude, similarly as above:

\gesperrt{$\delta$) If $n$ is not $= 6$, then the $n-2$ figures
$2, 3, \ldots, n-1$ are transitively connected by $Q_{0,1}$.}

For if they were not, then by $\alpha$):

\[
\begin{aligned}
&\text{either } \Pi(n-3) \text{ divisible by } g_1,\\
&\text{or } \Pi(n-4)\,\Pi(2) \geqq g_1
\end{aligned}
\]

would have to hold; hence by (10):

\[
\begin{aligned}
&\text{either } \frac{n}{n-2} \text{ an integer, hence } n \geqq 2n - 4,\\
&\text{or } n^2 - 7n + 6 \leqq 0.
\end{aligned}
\]

The first is not possible when $n > 4$, the second when $n > 6$.

Now if $n > 6$, the following also holds:

\gesperrt{$\varepsilon$) The group $Q_{0,1,2}$, which leaves the three figures 0, 1, 2
unchanged, has the degree}

\[ g_2 = \frac{g_1}{n-2} = \frac{\Pi(n-3)}{n}, \]

\gesperrt{and it is not possible that yet a fourth figure 3 should remain unchanged
under the whole group $Q_{0,1,2}$.}
